\section{Original whole twists on arbitrary schemes}

\begin{definition}[The actual whole original twist comparison]
\label{def:general-generated-line-actual-original-whole-twist-iso}
\uses{thm:general-generated-tuple-original-map-evaluation,def:fc-generated-line-whole-projective-map,thm:fc-generated-line-whole-projective-map-original-structure,thm:fc-coordinate-section-own-chart-frame,thm:fc-coordinate-section-chart-ratio,thm:fc-section-preserving-global-gluing}
\lean{EnvelopeProperty.GeneratedLineArbitrarySchemeTwist.projectiveTwistIso}
\leanok
Let $R$ be any commutative ring, $X$ the original arbitrary scheme,
$p:X\to\operatorname{Spec}R$ its original morphism, $L$ its original
invertible module, and $G$ its actual generating sections with a finite
index type. For every original coordinate pole $k$, construct an
isomorphism between the whole pullback of the original $O(1)$ along
the existing original generated-line morphism and the same original $L$.
There is no integrality, reducedness, generic point, desired morphism
or line-comparison certificate in the inputs.
\end{definition}
\begin{proof}
Reuse the existing whole generated-line morphism, which is already
defined over arbitrary schemes and has its actual original base
triangle. Reuse its cover by actual affine neighborhoods where one
original generator is a frame. On each such neighborhood the original
coefficient tuple spans the unit ideal and its pivot is exactly one.
The proved arbitrary-ring original tuple chart formula identifies the
same original affine map with its true projective coordinate factor.
The existing homogeneous-away section comparison identifies the
pullbacks of all original homogeneous coefficients with these actual
coefficients. The original projective pivot section is invertible on
its own chart, and all other sections have the same original coefficient
relative to it. Pull these facts through the actual chart factor and
the actual whole-map restriction triangle. Compare the inverse pulled
pivot with the original generator frame. The native scalar relations
show that this local isomorphism preserves every original generator.
The existing section-preserving gluing theorem derives overlap
compatibility from those section identities and the invertible original
pivot. It constructs the claimed whole original module isomorphism.
The old generic chart-cancellation identities and native pivot/gluing
core are reused; no affine fraction field is introduced.
\end{proof}

\begin{theorem}[All actual original sections are retained globally]
\label{thm:general-generated-line-actual-original-whole-twist-sections}
\uses{def:general-generated-line-actual-original-whole-twist-iso,thm:fc-section-preserving-global-gluing}
\lean{EnvelopeProperty.GeneratedLineArbitrarySchemeTwist.projectiveTwistIso_sectionHom}
\leanok
For every original generator index $j$, pull back the original global
homogeneous coordinate section indexed by the same original padded
generator coordinate. Composing it with the constructed whole
isomorphism gives exactly the same original global generator of $L$,
as a module homomorphism on the entire original $X$.
\end{theorem}
\begin{proof}
On every original affine generator frame, the coordinate section and
the original generator have the same actual original coefficient.
Cancellation of the actual pivot proves the local equality used in
the construction. Apply the global section conclusion of the existing
section-preserving gluing theorem to precisely this cover, these local
comparisons and these same original global sections. No additional
matching hypothesis or generic-stalk equality is supplied.
\end{proof}
